\subsection{20.1: a totally 2-closed insoluble group with nontrivial Fitting subgroup}
\label{cand:20.1}
The example is \(J_1\times C_{13}\).
The imported input is that \(J_1\) is totally \(2\)-closed
\cite[Theorem~1.1 and Proposition~4.8]{totally-two-closed}.
We give the product argument for arbitrary faithful actions, including
intransitive ones.

For \(X\le\operatorname{Sym}(\Omega)\), its \(2\)-closure consists
exactly of permutations \(s\) such that each ordered pair
\((\alpha,\beta)\) has an \(x\in X\) agreeing with \(s\) at both points.
In particular \(s\) preserves every point orbit.
For a finite group, total closure for finite faithful sets implies
the same for arbitrary sets: choose a finite invariant faithful subset
\(\Omega_0\), using one moved point for each nonidentity group element.
The restriction of \(s\) agrees with a unique group element there.
Applying finite closure to \(\Omega_0\cup\alpha^X\) for every \(\alpha\)
forces the same element everywhere.

\begin{lemma}
If finite groups \(A,B\) are totally \(2\)-closed and have coprime
orders, then \(A\times B\) is totally \(2\)-closed.
\end{lemma}
\begin{proof}
In any faithful action, \(\alpha^A\cap\alpha^B=\{\alpha\}\).
Indeed, \(\alpha^a=\alpha^b\) makes \(ab^{-1}\) fix \(\alpha\).
Since \(a,b\) commute and have coprime orders, the Chinese remainder
theorem expresses \(a\) as a power of \(ab^{-1}\), so \(a\) fixes
\(\alpha\).

Let \(\Delta_B\) be the set of \(B\)-orbits. The induced \(A\)-action
is faithful: an element fixing every such orbit moves each \(\alpha\)
within \(\alpha^A\cap\alpha^B\), so fixes every point.
Similarly \(B\) acts faithfully on \(\Delta_A\).
If \(s\) is in the \(2\)-closure of \(A\times B\), it preserves both
orbit equivalence relations, since each is a union of ordered-pair
orbits. For two \(B\)-orbits, choose representatives and apply the
pair criterion; the witnessing element \(ab\) induces the \(A\)-element
\(a\) on the pair of blocks.
Thus \(s\) induces an element of the \(2\)-closure of \(A\) on
\(\Delta_B\), and total closure gives one \(a_0\in A\) for all blocks.
Similarly one \(b_0\in B\) gives its action on \(\Delta_A\).
The permutation \(s(a_0b_0)^{-1}\) fixes every orbit of both factors
setwise, hence every singleton intersection, and is the identity.
\end{proof}

A prime cyclic group is totally \(2\)-closed: in a faithful action
choose a point with trivial stabilizer; pairing it with every other
point forces one common group element to realize a \(2\)-closure
permutation.
Now
\[
 |J_1|=175560=2^3\cdot3\cdot5\cdot7\cdot11\cdot19
\]
is coprime to \(13\), so the lemma proves the claim about closure.
The group is insoluble because it has the simple quotient \(J_1\).
Every normal nilpotent subgroup projects trivially to \(J_1\),
whereas the central \(C_{13}\) is nilpotent. Thus
\[
 F(J_1\times C_{13})=C_{13}\ne1.
\]
The supplied question imposes no direct-indecomposability restriction.
Prior related results include the nilpotent product case of
Chen--Ren and coprime transitive-action decompositions
\cite{chen-ren,abdollahi-arezoomand-tracey}; the lemma above handles
all faithful actions of the stated factors.
The \(J_1\) theorem retains its classification and computational
dependencies. See \artifact{research/20.1-proof.md}.
